% !TeX spellcheck = ru_RU
% !TeX root = vkr.tex

% Данные автора и руководителя перенесены из context/vkr.pdf.
\filltitle{ru}{
    chair              = {Кафедра системного программирования},
    group              = {23.Б-10},
    title              = {Реализация пакета параметрического оценивания для библиотеки pysatl-core},
    type               = {bachelor},
    kind               = {solution},
    author             = {Романюк Артем Дмитриевич},
    supervisorPosition = {доцент кафедры системного программирования, к.~ф.-м.~н.},
    supervisor         = {Гориховский~В.~И.},
    year               = {2026},
}

% Для постановки задачи используется краткий титульный лист.
% Коды направления и программы, отсутствующие в прошлой курсовой,
% не заполняются предположениями из шаблона.
\makeatletter
\renewcommand{\maketitleru}{
    \begin{titlepage}
        \newgeometry{top=20mm,bottom=20mm,left=20mm,right=20mm,nohead,nofoot}
        \thispagestyle{empty}
        \begin{center}
            {\large Санкт-Петербургский государственный университет\par}
            \vspace{1em}
            {\large \my@title@chair@ru\par}
            \vspace{1em}
            {\large Группа \my@title@group@ru\par}
            \vfill
            {\LARGE\bfseries \my@title@title@ru\par}
            \vspace{1.5em}
            {\Large \my@title@author@ru\par}
            \vspace{1em}
            {\large Выпускная квалификационная работа\par}
            \vspace{0.5em}
        \end{center}
        \vspace{3em}
        \begin{flushright}
            Научный руководитель:\\
            \my@title@supervisorPosition@ru\\
            \my@title@supervisor@ru
        \end{flushright}
        \vfill
        \begin{center}
            Санкт-Петербург\\
            \my@title@year@ru
        \end{center}
    \end{titlepage}
}
\makeatother
